\chapter{The annihilator and the count}

\begin{definition}[Rank hypothesis]
\label{def:rank-hyp}
\lean{ChabautyColeman.HypJac.rank}
\leanok
\uses{def:hypjac}
$\operatorname{rk} \J{\QQ} := \dim_{\QQ} \big(\J{\QQ} \otimes_{\ZZ} \QQ\big)$; the
standing hypothesis of the main theorem is $\operatorname{rk} \J{\QQ} \le r$ for a natural
number $r < g$. No finite generation is assumed: the dimension is a cardinal. The
qualitative form $\operatorname{rk} \J{\QQ} < g$ is equivalent to the existence of such an
$r$, a cardinal below a natural number being itself a natural number.
\emph{Formalization:} the Jacobian is the class group, written multiplicatively, so the
rationalization is $\QQ \otimes_{\ZZ}$ of the additivization; the arithmetic base field
is an arbitrary subfield $K_0$ of the local field, not just $\QQ$.
\end{definition}

\begin{lemma}[The evaluation map factors through $\otimes\,\QQ$]
\label{lem:eval-factors}
\lean{ChabautyColeman.lambdaEval, ChabautyColeman.lambdaEvalRat}
\leanok
\uses{thm:glue, def:basechange, def:rank-hyp}
The bilinear map $(\omega, c) \mapsto \lambda_\omega(c_{\Qp})$, for $c \in \J{\QQ}$
mapped to $\J{\Qp}$ by the homomorphism of Definition~\ref{def:basechange}, induces a
$\QQ$-linear map
$\J{\QQ} \otimes \QQ \to \operatorname{Hom}_{\Qp}(\RegDiff, \Qp)$.
\end{lemma}
\begin{proof}
\leanok
Homomorphisms into the torsion-free uniquely divisible group $\Qp$ factor through
$\otimes\,\QQ$. \emph{Formalization.} The base change is built into $\lambda$
(Definition~\ref{def:lambda}), and the differential is represented by its coordinate
vector in the $\omega'$-basis; the statement is proved for any pairing of an abstract
commutative group that is additive in the class and linear in the coordinates, with
these three properties as hypotheses, and instantiated at $\lambda$.
\end{proof}

\begin{theorem}[Annihilating differential]
\label{thm:annihilator}
\lean{ChabautyColeman.exists_annihilator_submodule,
  ChabautyColeman.HypJac.rank_subgroup_le,
  ChabautyColeman.sum_smul_tracePoly_ne_zero,
  Module.finrank_le_finrank_iInf_ker_add_finrank}
\leanok
\uses{lem:eval-factors, def:rank-hyp, def:regdiff}
If $\operatorname{rk}\J{\QQ} \le r$, the differentials $\omega \in \RegDiff$ with
$\lambda_\omega(c) = 0$ for every $c \in J^1(\QQ)$ form a subspace of codimension at most
$r$; for any such $\omega$, in particular,
$\lambda_\omega\big(\aj(P) - \aj(P_0)\big) = 0$ for all $P, P_0 \in C(\QQ)$ in one
residue disk. For $r < g$ the subspace is nonzero.
\end{theorem}
\begin{proof}
\leanok
$\operatorname{rk} J^1(\QQ) \le \operatorname{rk}\J{\QQ}$, since $\otimes\,\QQ$
preserves the inclusion $J^1 \subseteq \J{\QQ}$ by flatness of $\QQ$ over $\ZZ$. So the
image of Lemma~\ref{lem:eval-factors} spans a $\Qp$-subspace of dimension
$\le \operatorname{rk}\J{\QQ} \le r$, and the annihilator of that span in $\RegDiff$ has
codimension at most $r$. The dimension count is isolated as an abstract statement: for an
$R$-linear map $Q \to \operatorname{Hom}_L(V, L)$, with $R$ a subfield of $L$, the
subspace of $V$ annihilated by the whole image has codimension at most $\dim_R Q$,
because the $L$-span of the image is generated by the images of an $R$-basis of the
source, so the annihilator is the common kernel of $\dim_R Q$ functionals. The
differentials are produced as their coordinate vectors in the $\omega'$-basis; that a
nonzero vector represents a nonzero element of $\RegDiff$ is the linear independence of
the trace-form numerators (Theorem~\ref{thm:omega-prime-global}(c)).
\end{proof}

\begin{lemma}[Integral normalization]
\label{lem:normalize}
\lean{ChabautyColeman.exists_annihilator_residue,
  Polynomial.exists_submodule_map_residue,
  IsLocalRing.linearIndependent_map_of_saturated,
  ChabautyColeman.exists_annihilator_tracePoly}
\leanok
\uses{thm:annihilator, def:regdiff, def:hypcurve}
The reductions $\bar h$ of the integral numerators of the annihilating differentials form
a subspace $V$ of the polynomials of degree $< g$ over $\Fpbar$, again of codimension at
most $r$, and every element of $V$ is the reduction of such a numerator. In particular a
nonzero element of $V$ comes from an $\omega \in \RegDiff_{\ZZ}$ with $\bar\omega \ne 0$
on $\bar C_{\Fpbar}$.
\end{lemma}
\begin{proof}
\leanok
The annihilating differentials with integral numerator form a lattice in the annihilating
subspace of Theorem~\ref{thm:annihilator}: a $\Zp$-submodule of the numerators of degree
$< g$, free of rank equal to the dimension of that subspace, and saturated, since
$\pi^{-1}h$ is again integral and annihilating whenever $\pi \mid h$. Reduction is
therefore injective on a basis of the lattice, and the span of the reduced basis is the
required $V$; an element of $V$ lifts by lifting the coefficients of that basis, scaling
does not affect the vanishing of $\lambda_\omega$, which is linear.
\emph{Formalization.} The differentials are handled through their numerators $h$ (the
polynomials with $\omega = h\,dx/y$, $\deg h < g$), which is what the order computations
of Lemma~\ref{lem:ord-match} and Corollary~\ref{cor:sum-ord-space} consume; the lattice
statement \texttt{exists\_submodule\_map\_residue} is proved for an arbitrary subspace of
the polynomials of bounded degree over the fraction field of a discrete valuation ring.
The composed form used by the main theorem is \texttt{exists\_annihilator\_tracePoly}:
Theorem~\ref{thm:annihilator} instantiated in the trace-form basis, giving for each
$\bar h \in V$ coordinates $c$ \emph{and} an integral numerator $h_0$ with reduction
$\bar h$ and degree $< g$ whose image is
$\sum_i c_i\,\bigl(\text{trace-form numerator}\bigr)_i$ --- a statement about $\tau$
alone, independent of any residue disk, so every residue disk may pick its own
differential out of $V$.
\end{proof}

\begin{lemma}[Disk decomposition of the rational points]
\label{lem:disk-decomp}
\lean{ChabautyColeman.ratPoints_disk_decomp,
  ChabautyColeman.encard_le_of_fiber_encard_le,
  ChabautyColeman.exists_finset_redPoint_rational}
\uses{lem:two-charts, def:disk}
\leanok
$C(\QQ) \subseteq C(\Qp) = \bigsqcup_{\bar P \in \bar C(\Fp)} \big(C(\Qp) \cap
D_{\bar P}\big)$, and rational points reduce to $\Fp$-rational residue points.
\end{lemma}
\begin{proof}
\leanok
Lemma~\ref{lem:two-charts} over $O = \Zp$; coordinates in $\Zp$ (or $t \in p\Zp$)
reduce into $\Fp$. \emph{Formalization.} The decomposition is definitional --- the
residue disk \emph{is} the fibre of the reduction map --- so what the count uses is the
fibrewise bound: a set of points reducing into a finite set $T$ of residue points, with
at most $1 + m_{\bar P}$ points on the disk of each $\bar P \in T$, has at most
$\#T + \sum_{\bar P \in T} m_{\bar P}$ elements. Cardinalities are extended naturals,
so finiteness is not assumed but obtained. The arithmetic half --- that reductions of
$K_0$-rational points land in a set of size at most the point count of the reduced
curve over the residue field of an arithmetic valuation ring $O_0 \subseteq K_0$
compatible with $O$ --- is \texttt{exists\_finset\_redPoint\_rational}.
\end{proof}

\begin{theorem}[Per-disk bound]
\label{thm:disk-bound}
\lean{ChabautyColeman.encard_le_of_lambda_eq_zero, ChabautyColeman.disk_count_le,
  ChabautyColeman.disk_bound_generic, ChabautyColeman.disk_bound_weierstrass,
  ChabautyColeman.disk_bound_inf,
  ChabautyColeman.encard_inter_residueDisk_le,
  ChabautyColeman.encard_inter_residueDisk_le_ordPoint,
  ChabautyColeman.encard_inter_residueDisk_le_ordPoint_tracePoly,
  ChabautyColeman.encard_inter_residueDisk_le_ordPoint_tracePoly_weierstrass,
  ChabautyColeman.encard_inter_residueDisk_le_ordPoint_tracePoly_inf,
  ChabautyColeman.map_omegaExpansionWeierstrass_eq_sum_smul_omegaPrime_tracePoly,
  ChabautyColeman.map_omegaExpansionInf_eq_sum_smul_omegaPrime_tracePoly,
  ChabautyColeman.PSeries.encard_le_of_evalT_const,
  ChabautyColeman.PSeries.hasSum_algebraMap_evalT,
  ChabautyColeman.PSeries.evalT_sub_C,
  ChabautyColeman.PSeries.exists_substituted_antideriv,
  IsLocalRing.natCast_dvd_pow_self,
  ChabautyColeman.PSeries.algebraMap_coeff_eq_coeff_antideriv}
\uses{thm:glue, thm:annihilator, lem:normalize, lem:ord-match, thm:antideriv-zeros,
  def:disk-antider, cor:sum-ord-space, def:hypcurve}
\leanok
Let $\bar P \in \bar C(\Fp)$ and $m := m_{\bar P} = \ord_{\bar P}(\bar\omega)$. Then
$\#\big(C(\QQ) \cap D_{\bar P}\big) \le 1 + m$.
\end{theorem}
\begin{proof}
\leanok
If the intersection is empty, we are done; else fix $P_0 \in C(\QQ) \cap D_{\bar P}$.
For any other rational $P$ in the disk, Theorems~\ref{thm:annihilator}
and~\ref{thm:glue} give
$0 = \lambda_\omega(\aj P) = \lambda_\omega(\aj P_0) + G^{\omega}_D(z(P)) -
G^{\omega}_D(z(P_0))$, so $z(P)$ is a zero in $p\Zp$ of the single function
$G^{\omega}_D(z) + c_0$. The derivative series of $G$ has reduction of order exactly
$m$ (Lemmas~\ref{lem:normalize} and~\ref{lem:ord-match}); choosing for $\omega$ the
differential of smallest order at $\bar P$ gives $m \le 2r$
(Corollary~\ref{cor:sum-ord-space}), and $p > 2r+2$ gives $p > m + 2$; apply
Theorem~\ref{thm:antideriv-zeros}.

\emph{Formalization.} The analytic core is: for points of one disk with distinct
parameters in $\mathfrak m$ at which the integral takes one and the same value,
$\#S \le 1 + m$; a constant value is as good as a zero (subtracting it changes only the
constant coefficient), and the value delivered by Theorem~\ref{thm:glue} --- a
convergent sum over $\mathrm{Frac}(O)$ --- is identified coefficientwise with the value
of the substituted series $H(u) = G(\pi u)$ over $O$ on which
Theorem~\ref{thm:antideriv-zeros} runs; integrality of $H$ is the divisibility of
$\pi^{k+1}$ by $k+1$ in $O$, which is where the bound on the ramification is spent. Rationality of the points enters only through the coordinate
identities matching each $P \in S$ with the parametrized disk point at its parameter:
these give distinctness of the parameters and, through Theorem~\ref{thm:glue} and
Theorem~\ref{thm:annihilator}, the constancy of the integral. The identification of the
disk expansion of $\omega = \sum_i c_i\,\omega'_i$ with the integrand --- the remaining
hypothesis --- is supplied on each of the three chart types by the per-disk identities
of Theorem~\ref{thm:omega-prime-global}(c), fed with the annihilating differential
normalized in the trace-form basis (Lemma~\ref{lem:normalize}); the statements
\texttt{disk\_bound\_generic}, \texttt{disk\_bound\_weierstrass} and
\texttt{disk\_bound\_inf} establish the bound at each chart type from the data of
Lemma~\ref{lem:tau-exists} alone, with the centre lift, the auxiliary point and all
B\'ezout data constructed inside. All of these statements, like Lemma~\ref{lem:normalize}
and the counting step of Theorem~\ref{thm:main-count}, are formulated for an odd-degree
hyperelliptic curve \texttt{C : HyperellipticCurve O} (Definition~\ref{def:hypcurve}) and a
curve \texttt{C₀} over $K_0$: the curve supplies monicity, the degree, oddness and
squarefreeness over $\operatorname{Frac} O$, over $K_0$ and over the residue field through
its instances, and the invertibility of $2$ in $O$ is the instance \texttt{[Invertible (2 : O)]}.
\end{proof}

\begin{theorem}[Chabauty--Coleman, local count]
\label{thm:main-count}
\lean{ChabautyColeman.chabauty_coleman_odd_hyperelliptic,
  ChabautyColeman.ratPoints_encard_le,
  ChabautyColeman.main_affine, ChabautyColeman.ratPoints_finite_of_rank_le,
  ChabautyColeman.CurvePoint.encard_univ,
  Polynomial.separable_map_of_separable_residue}
\uses{thm:disk-bound, lem:disk-decomp, cor:sum-ord-space, def:rank-hyp, def:hypcurve}
\leanok
Under the standing hypotheses, let $K_0$ be a subfield of $\operatorname{Frac} O$ of
characteristic $0$, and let $S$ be a set of points of $C$ with $K_0$-rational coordinates
whose reductions lie in a finite set $T$ of points of $\bar C$. If
$\operatorname{rk}\J{K_0} \le r$ with $r < g$ and $p > 2r + e + 1$, then
\[ \#S\ \le\ \#T\ +\ 2r. \]
The statement over $\QQ$ (Theorem~\ref{thm:main}) follows via the arithmetic form
(Theorem~\ref{thm:main-arith}) and its number-field instantiation
(Theorem~\ref{thm:main-nf}).
\end{theorem}
\begin{proof}
\leanok
Sum Theorem~\ref{thm:disk-bound} over the disks of Lemma~\ref{lem:disk-decomp}, each disk
using the annihilating differential of smallest order at its centre:
$\#S \le \sum_{\bar P \in T} \bigl(1 + n(V, \bar P)\bigr)
\le \#T + 2r$ by Corollary~\ref{cor:sum-ord-space}.

\emph{Formalization.} The full assembly is
\texttt{chabauty\_coleman\_odd\_hyperelliptic}, stated for an odd-degree hyperelliptic curve
\texttt{C : HyperellipticCurve O} (its unit discriminant is the good reduction) and a curve
\texttt{C₀} over $K_0$ agreeing with it over $\operatorname{Frac} O$; the polynomial-level
hypotheses of the inner layers are supplied by the curves. The proof constructs the class $\tau$ avoiding the
$x$-coordinates of $T$ (Lemma~\ref{lem:tau-exists} via Lemma~\ref{lem:tau-generic},
unconditionally for every genus $\ge 1$), produces the space $V$ of reduced numerators of
the annihilating differentials in the trace-form basis (Theorem~\ref{thm:annihilator},
Lemma~\ref{lem:normalize}), picks in each residue disk an element of $V$ of minimal
order there, classifies the residue point as generic affine, Weierstrass or $\infty$,
invokes the matching per-disk bound of Theorem~\ref{thm:disk-bound}, and sums. No genus
lower bound is assumed: the rank hypothesis forces $g \ge 1$, and genus
$1$ passes through Lemma~\ref{lem:tau-generic} since $\bar f$ is squarefree. The
hypotheses of the formal statement are exactly: $O$ a complete discrete valuation ring
of characteristic $0$ with algebraically closed residue field of characteristic
$\neq 2$ --- carrying a uniformizer $\pi$ and a bound $e$ on the absolute ramification
index, with no unramifiedness assumed --- $f$ monic with squarefree reduction of odd
degree, and the rank hypothesis
--- all other instances (oddness and squarefreeness over the fraction field and the
base field, adic completeness) are derived inside; for the squarefreeness this is the
statement that separability of the reduction of a monic polynomial over a local ring
lifts through any ring homomorphism, by the resultant
(\texttt{separable\_map\_of\_separable\_residue}). The affine restatement --- both sides
counting the projective sets with the point at infinity made explicit, of which there is
exactly one on each side by the odd degree --- is \texttt{main\_affine} (via
\texttt{CurvePoint.encard\_univ}), and the qualitative finiteness of any set of
base-rational points is \texttt{ratPoints\_finite\_of\_rank\_le}: the counts are extended
naturals, so the bound caps the cardinality by a natural number.
\end{proof}

\begin{theorem}[Chabauty--Coleman over an arithmetic base]
\label{thm:main-arith}
\lean{ChabautyColeman.chabauty_coleman_odd_hyperelliptic_arith,
  ChabautyColeman.main_affine_arith, ChabautyColeman.ratPoints_finite_arith,
  IsLocalRing.IsUniformizer}
\uses{thm:main-count, lem:disk-decomp, thm:antideriv-zeros}
\leanok
Let $O$ be a complete discrete valuation ring of characteristic $0$ with algebraically
closed residue field of characteristic $\ne 2$, uniformizer $\pi$, residue characteristic
$p$ and absolute ramification index at most $e$ (that is, $v(p) \le e$). Let $K$ be a
subfield of $\operatorname{Frac} O$ of characteristic $0$ and $O_0$ a local ring mapping
compatibly into $O$ and into $K$, with finite residue field $\Fq$, such that every element
of $K$ of valuation $\le 1$ comes from $O_0$. Let $C \colon y^2 = f(x)$ be a hyperelliptic
curve of genus $g$ over $O$ with $f$ monic of odd degree $2g+1$ and coefficients in $O_0$
(so its reduction over the residue field of $O$ is squarefree). If
$\operatorname{rk} J(K) \le r$ with $r < g$ and
\[ p \;>\; 2r + e + 1, \]
then
\[ \#C(K)\ \le\ \#\bar C(\Fq)\ +\ 2r. \]

For a number field $K$ and a prime $\mathfrak p$ of good reduction the
hypotheses hold with $O$ the ring of integers of the completion of the maximal unramified
extension of $K_{\mathfrak p}$, $O_0$ the localization of $O_K$ at $\mathfrak p$,
$q = \#(O_K/\mathfrak p)$ and $e = e(\mathfrak p \mid p)$, so ramified primes are admitted:
this is Corollary~6.7 of Stoll, \emph{loc.\ cit.}; that instantiation is
Theorem~\ref{thm:main-nf}, and its case $K = \QQ$ is Theorem~\ref{thm:main}.
\end{theorem}
\begin{proof}
\leanok
Specialize Theorem~\ref{thm:main-count}: the base field enters only
through the rationality of the coordinates of the points being counted, and the
ramification only through the substitution $z = \pi u$ of
Theorem~\ref{thm:antideriv-zeros}, which asks $p > m + e + 1$ for the order $m$ counted on
the disk; since $m \le 2r$ at every disk (Corollary~\ref{cor:sum-ord-space}), the condition
$p > 2r + e + 1$ suffices. The passage from the residue points over the algebraically closed
residue field of $O$ to the $\Fq$-points of $\bar C$ is the arithmetic half of
Lemma~\ref{lem:disk-decomp}.

\emph{Formalization.} The arithmetic base is carried as a pack: a local ring $O_0$ with
finite residue field together with ring maps $O_0 \to O$ (local) and $O_0 \to K$ that agree
in $\operatorname{Frac} O$, the hypothesis that $O_0$ covers $\{x \in K : v(x) \le 1\}$, and
a polynomial over $O_0$ mapping to $f$; the ramification data is
\texttt{IsLocalRing.IsUniformizer $\pi$ $p$ $e$} ($\mathfrak m = (\pi)$, $p \in \mathfrak m$,
$p \mid \pi^e$). The affine forms \texttt{main\_affine\_arith} and
\texttt{ratPoints\_finite\_arith} state
$\#\{(x,y) \in K^2 : y^2 = f(x)\} + 1 \le \#\{(\bar x,\bar y) \in \mathbb F_q^2 :
\bar y^2 = \bar f(\bar x)\} + 1 + 2r$ and the finiteness of the affine $K$-solution set.
The pack is instantiated in Theorem~\ref{thm:main-nf}, at the local factor of
Chapter~\ref{ch:numberfield} with $O_0$ the localization at $\mathfrak p$; the statements
over $\QQ$ follow from that instantiation (Theorem~\ref{thm:main},
Corollary~\ref{cor:finiteness}).
\end{proof}

\begin{theorem}[Chabauty--Coleman over $\QQ$]
\label{thm:main}
\lean{ChabautyColeman.main_affine_rat_of_rank_le, ChabautyColeman.main_affine_int_of_rank_le,
  ChabautyColeman.main_affine_int, ChabautyColeman.separable_of_squarefree_twist}
\uses{thm:main-nf}
\leanok
Let $C \colon y^2 = f(x)$ be a hyperelliptic curve of genus $g$ over $\QQ$ with $f$ monic of
odd degree $2g+1$, and $p$ a prime such that $f$ has $p$-integral coefficients and
$\bar f \in \Fp[x]$ is squarefree. If $\operatorname{rk} \J{\QQ} \le r$ with $r < g$ and
$p > 2r + 2$, then
\[ \#C(\QQ)\ \le\ \#\bar C(\Fp)\ +\ 2r. \]
\end{theorem}
\begin{proof}
\leanok
This is the case $A = \ZZ$, $K = \QQ$, $\mathfrak p = (p)$ of
Theorem~\ref{thm:main-nf}: the residue field of $(p)$ is $\Fp$, and the ramification
index is $1$ (it is at most $[\QQ:\QQ]$), so the prime condition $p > 2r + e + 1$ reads
$p > 2r + 2$.

\emph{Formalization.} The Lean statement is the affine form --- both sides count the
projective sets with the point at infinity made explicit, of which there is exactly one
on each side by the odd degree --- with the reduction count over $\mathrm{ZMod}\ p$;
$p$-integrality is certified by $d f = F \in \ZZ[x]$ with $p \nmid d$, and the reduced
count is taken on the model $\bar d\, y^2 = \bar F(x)$, to which the specialization
transports squarefreeness and the count along $\ZZ_{(p)}/(p) \cong \mathrm{ZMod}\ p$
(\texttt{main\_affine\_rat\_of\_rank\_le}, stated for \texttt{C : HyperellipticCurve ℚ}; the
comparator builds the curve from $f$ by \texttt{HyperellipticCurve.ofMonicOdd}, the squarefree
reduction making $f$ separable, \texttt{separable\_of\_squarefree\_twist}). The version for
$f \in \ZZ[x]$ is its case $d = 1$ (\texttt{main\_affine\_int\_of\_rank\_le}), whose case $r = g - 1$ is Coleman's
bound \texttt{main\_affine\_int}, where the prime condition reads $p \ge \deg f$.
\end{proof}

\begin{corollary}[Finiteness]
\label{cor:finiteness}
\lean{ChabautyColeman.NumberFieldBase.ratPoints_finite_numberField}
\uses{cor:finiteness-nf}
\leanok
If $f \in \QQ[x]$ is monic, squarefree and of odd degree $2g+1$, and
$\operatorname{rk} \J{\QQ} < g$, then $C(\QQ)$ is finite.
\end{corollary}
\begin{proof}
\leanok
This is the case $A = \ZZ$, $K = \QQ$ of Corollary~\ref{cor:finiteness-nf}; it is not
restated separately in Lean.
\end{proof}

\begin{remark}[Genus one, rank zero]
\label{rem:genus-one}
\uses{thm:main}
Theorem~\ref{thm:main} carries no genus lower bound: the rank hypothesis forces
$g \ge 1$, and the genus-one case is covered unconditionally
(Lemma~\ref{lem:fold-witness-genus-one}). At $g = 1$ the rank hypothesis reads
$\operatorname{rk} = 0$, so $r = 0$, and the bound is $\#E(\QQ) \le \#\bar E(\Fp)$ for a
rank-zero elliptic curve $y^2 = \text{cubic}$ at any odd prime of good reduction --- the
prime condition $p > 2r + 2$ is exactly oddness.
\end{remark}
